% !PRINTED 449
\weberpage{449}
%<<CONT
then by 5. $N(p) = p^{2h}$, and consequently
\[
2h = r\varepsilon.
\]

The number of prime factors into which $p$ splits in the field $\mathfrak{K}$
is therefore equal to $2h/\varepsilon$.

If further $(D, p) = -1$, then the prime factors of $p$ are of the second degree, and their number is therefore $= h$.

9. We draw one further interesting conclusion from these results. Let $(D, p) = +1$ and let $\varepsilon$ be the index of $p$. Then $4p^{\varepsilon}$ is representable by the form $x^2 - Dy^2$, and
\weq{23}{p^{\varepsilon} = \mu\mu',}
if for brevity
\weq{24}{\mu = \frac{x + y\sqrt{D}}{2}, \quad \mu' = \frac{x - y\sqrt{D}}{2}}
is set. We now know that $p$ splits into nothing but distinct prime ideals of degree $\varepsilon$. Since $p$ is odd, and $x, y$ have at most the common divisor 2, $\mu, \mu'$ have no common prime factor in $\mathfrak{K}$. To every prime ideal $\mathfrak{P}$ that divides $\mu$ there corresponds a prime ideal $\mathfrak{P}'$ distinct from it (the conjugate one), which arises from $\mathfrak{P}$ by putting, for all numbers of $\mathfrak{P}$, the conjugate imaginary numbers. Thus $\mu$ is divisible not only by $\mathfrak{P}$, but by $\mathfrak{P}^{\varepsilon}$.

According to \S\,121 we have
\weq{25}{\left(\frac{\eta(\omega_1)}{\eta(\omega)}\right)^2 \left(\frac{\eta(\omega_2)}{\eta(\omega_1)}\right)^2 \cdots \left(\frac{\eta(\omega_{\varepsilon})}{\eta(\omega_{\varepsilon-1})}\right)^2 = \varrho\mu,}
and by \S\,72, 4 the factors $P^2$ on the left side of this expression satisfy an equation $P^2\varphi(P^2, \gamma_2, \gamma_3) = p$, in which $\varphi(P^2, \gamma_2, \gamma_3)$ is an integral algebraic number.

From this it follows, however, that $P^2$ cannot be divisible by a higher than the first power of $\mathfrak{P}$, while $\mu$ is nevertheless divisible by $\mathfrak{P}^{\varepsilon}$. We therefore conclude from (25) that \gesperrt{each} of the numbers $P^2$ must be divisible by the first power of $\mathfrak{P}$, and hence that all these numbers \gesperrt{are associated}. If, then, $\varrho$ denotes any algebraic unit, we have
\[
\frac{x + y\sqrt{-m}}{2} = \varrho P^{2\varepsilon},
\]
so $\mu$ is really represented as the $\varepsilon$th power of a number existing in the field $\mathfrak{K}$, and $P^2$ itself is a representation of a prime ideal in the field $\Omega$ by a \gesperrt{principal ideal} in the field $\mathfrak{K}$.
